\documentclass[leqno]{jsarticle}
\usepackage{amssymb}
\usepackage{amsthm}
\usepackage{amsmath}
\usepackage{comment}
	%可換図式用
		%\usepackage[all]{xy}
	%重くなるので使わいときはコメント化しておく。
%定理環境 thmはあまり使わずpropをなるべく使う。
%主張は基本的に証明の中で使い、そのため番号を付けない。
%注意環境を付けてみた。
\newtheorem{thm}{定理}
\newtheorem{prop}[thm]{命題}
\newtheorem{cor}[thm]{系}
\newtheorem{lem}[thm]{補題}
\newtheorem{df}[thm]{定義}
\newtheorem{exam}[thm]{例}
\newtheorem{fact}[thm]{事実}
\newtheorem*{claim}{主張}
\newtheorem*{cau}{注意}
\renewcommand{\proofname}{{\bf 証明}}
%矢印たち。
%\toは\rightarrowと同じ。\getsは\leftarrowと同じ。同値は\iff
%上向き矢はuparrow逆はdownarrow
\newcommand{\To}{\Longrightarrow}
\newcommand{\Gets}{\Longleftarrow}
%黒板太字
\newcommand{\nn}{\mathbb{N}}
\newcommand{\zz}{\mathbb{Z}}
\newcommand{\qq}{\mathbb{Q}}
\newcommand{\rr}{\mathbb{R}}
\newcommand{\cc}{\mathbb{C}}
%無限大
\newcommand{\mgn}{\infty}
%ギリシャン
\newcommand{\dpst}{\displaystyle}
\newcommand{\ep}{\varepsilon}
\newcommand{\al}{\alpha}
\newcommand{\be}{\beta}
\newcommand{\de}{\delta}
\newcommand{\la}{\lambda}
\newcommand{\La}{\Lambda}
\newcommand{\ga}{\gamma}
\newcommand{\lil}{\la\in \La}
%商集合
\newcommand{\qt}[2]{#1/\! #2}
%enumerate環境
\renewcommand{\labelenumi}{{\rm (\arabic{enumi})}}
%以降alignじゃなくてenumerate を使え。
%なんか演算
\newcommand{\card}{\text{{\rm card}}}
\newcommand{\supp}{\text{{\rm supp}}}
\newcommand{\wsp}{\text{{\rm wspp}}}
\newcommand{\wspp}{\text{{\rm wspp}}}
%なんか演算２
\newcommand{\bd}{\text{{\rm Bdry}}}
\newcommand{\cl}{\text{{\rm CL}}}
\newcommand{\nai}{\text{{\rm Int}}}
%差集合は\setminusで統一する。等号の否定は\neq
%真部分集合は\subsetneqq　偏微分は\partial
%ディスプレイスタイル。\dfracでディスプレイ分数
%空集合は\Oを使う！！！！！！！！！！！！

\title{単位の分割}
\author{一色三良(concious77)}
\date{\today}
			\begin{document}
\maketitle
この文章では正規空間の様々な特徴付けを紹介する。ここでは正規空間とは$T_1$を仮定しない意味で用いる。

\begin{df}[細分]$\mathfrak{A},\mathfrak{B}を$集合$X$の部分集合の族とする。このとき$\mathfrak{A}$が$\mathfrak{B}$の細分であるとは
\[
\forall B\in \mathfrak{B}\ \exists A\in \mathfrak{A}\ A\subset B
\]
となることである。
\end{df}

\begin{df}[点有限族]集合$X$の部分集合の族$\mathfrak{A}$が点有限であるとは$X$の任意の点$x$に対して$x\in A$となる$A\in \mathfrak{A}$が有限個しかないことである。

\end{df}

\begin{df}[局所有限族]位相空間$X$に於いて$X$の部分集合の族$\mathfrak{A}$が局所有限であるとは$X$の各点$x$について$x$の近傍$N$が存在して
\[
N\cap A\neq \emptyset となるA\in \mathfrak{A}は有限個
\]
を充たすことである。
\end{df}

\begin{df}[開収縮]位相空間$X$の開被覆$\mathcal{U}$に対して$\{V_U\ |\ U\in \mathcal{U}\}$が開収縮であるとは$\{V_U\ |\ U\in \mathcal{U}\}$が$X$の開被覆であって、任意の$U\in \mathcal{U}$について
\[
CL(V_U)\subset U
\]
を充たすこと。
\end{df}

\begin{df}[閉収縮]位相空間$X$の開被覆$\mathcal{U}$に対して$\{F_U\ |\ U\in \mathcal{U}\}$が閉収縮であるとは$\{F_U\ |\ U\in \mathcal{U}\}$が$X$の閉被覆であって、任意の$U\in \mathcal{U}$について
\[
F_{U}\subset U
\]
を充たすこと。
\end{df}


\begin{df}[Weak supportとSupport]
関数$fX\to [0,1]$の{\bf 弱台}$\wspp (f)$を
\[
\wspp (f)=\{x\in X\ |\ f(x)>0\}
\]
と定義する。\par
関数$fX\to [0,1]$の{\bf 台}$\supp (f)$を
\[
\supp (f)=CL(\supp(f))=CL(\{x\in X\ |\ f(x)>0\})
\]

\end{df}

以下様々な種類の単位の分割を定義する。
%弱い単位の分割
\begin{df}[弱い単位の分割]位相空間$X$上の実数値連続関数の族$\{f_i\}_{i\in I}$が$X$の開被覆$\mathcal{U}$に従属する弱い単位の分割であるとは$\{\wspp f_i\}_{i\in I}$が$\mathcal{U}$の細分であり
\begin{align*}
&\sum_{i\in I} f_i=1
\end{align*}
を充たすこと
\end{df}
%強い単位の分割
\begin{df}[強い単位の分割]位相空間$X$上の実数値連続関数の族$\{f_i\}_{i\in I}$が$X$の開被覆$\mathcal{U}$に従属する強い単位の分割であるとは$\{\supp f_i\}_{i\in I}$が$\mathcal{U}$の細分であり
\begin{align*}
&\sum_{i\in I} f_i=1
\end{align*}
を充たすこと
\end{df}
%点有限な弱い単位の分割
\begin{df}[点有限な弱い単位の分割]位相空間$X$上の実数値連続関数の族$\{f_i\}_{i\in I}$が$X$の開被覆$\mathcal{U}$に従属する弱い単位の分割であるとは$\{\wspp f_i\}_{i\in I}$が$\mathcal{U}$の細分であり
\begin{align*}
&\sum_{i\in I} f_i=1\\
&\text{$\{\wspp f_i\}_{i\in I}$は点有限}
\end{align*}
を充たすこと
\end{df}

%点有限な強い単位の分割
\begin{df}[点有限な弱い単位の分割]位相空間$X$上の実数値連続関数の族$\{f_i\}_{i\in I}$が$X$の開被覆$\mathcal{U}$に従属す強い単位の分割であるとは$\{\supp f_i\}_{i\in I}$が$\mathcal{U}$の細分であり
\begin{align*}
&\sum_{i\in I} f_i=1\\
&\text{$\{\supp f_i\}_{i\in I}$は点有限}
\end{align*}
を充たすこと
\end{df}

%局所有限な弱い単位の分割
\begin{df}[局所有限な弱い単位の分割]位相空間$X$上の実数値連続関数の族$\{f_i\}_{i\in I}$が$X$の開被覆$\mathcal{U}$に従属する局所有限な弱い単位の分割であるとは$\{\wspp f_i\}_{i\in I}$が$\mathcal{U}$の細分であり
\begin{align*}
&\sum_{i\in I} f_i=1\\ 
&\text{$\{\wspp f_{i}\}_{i\in I}$が局所有限}
\end{align*}
を充たすこと
\end{df}

%局所有限な強い単位の分割
\begin{df}[局所有限な強い単位の分割]位相空間$X$上の実数値連続関数の族$\{f_i\}_{i\in I}$が$X$の開被覆$\mathcal{U}$に従属する局所有限な強い単位の分割であるとは$\{\supp f_i\}_{i\in I}$が$\mathcal{U}$の細分であり
\begin{align*}
&\sum_{i\in I} f_i=1\\ 
&\text{$\{\supp f_{i}\}_{i\in I}$が局所有限}
\end{align*}
を充たすこと
\end{df}

\begin{comment}
\begin{df}[広義の単位の分割]
関数族$\{f_{\la}\}_{\lil}$が$X$の開被覆$\mathcal{U}$に従属する単位の分割であるとは$\{\wsp(f_{\la})\}$が$\mathcal{U}$の細分になっており
\[
\sum_{\lil}f_{\la}\equiv 1
\]
を充たすときに言う。
\end{df}

\begin{df}[広義の局所有限な単位の分割]
関数族$\{f_{\la}\}_{\lil}$が$X$の開被覆$\mathcal{U}$に従属する単位の分割であるとは$\{\wsp(f_{\la})\}$が局所有限で$\mathcal{U}$の細分になっており
\[
\sum_{\lil}f_{\la}\equiv 1
\]
を充たすときに言う。
\end{df}
\end{comment}

\begin{lem}
$\{f_{\la}\}_{\lil}$が弱い単位の分割とすると、任意の$x\in X$及び任意の$\ep>0$に対して$\La$の有限集合$A$が存在して
\[
\sum_{\la\in A}f_{\la}(x)>1-\ep
\]
が成り立つ。このとき$x\in \{f_{b}>\ep\}$ならば$b\in A$が成り立つ。
\end{lem}

\begin{lem}$\{f_{\la}\}_{\lil}$が弱い単位の分割とし、
任意の$\de>0$について
\[
V_a^{\de}=\{f_a>\de\}
\]
と定義すると$\{V_{\la}^{\de}\}_{\lil}$は$X$で局所有限である。
\end{lem}
\begin{proof}
$p\in X$と$\de$について先の補題の有限集合$A$を取り
\[
N=\{\sum_{\la\in A}f_{\la}>1-\de\}
\]
を考えると$N\cap V_{a}^{\de}\neq \O$ならば先の補題から$a\in A$なので$N$は族の有限個の元としか交わらない。
\end{proof}

\begin{lem}
弱い単位の分割$\{f_{\la}\}_{\lil}$について
\[
F=\sup_{\lil}f_{\la}
\]
は連続である。
\end{lem}
\begin{proof}
$p\in X$を与える。このときある$a\in \La$について$f_a(p)>0$なので適当に$f_a(p)>\de$となる$\de>0$を取ると
$\{V_{\la}^{\de}\}_{\lil}$は先の補題から局所有限である。その局所有限性を担保する点$p$における近傍を$N$とし、有限集合$A$は
\[
N\cap V_{\la}^{\de}\neq \O\To \la \in A
\]
となるものとする。このとき$N$上で
\[
F=\max_{a\in A}f_{a}
\]
なので$F$は$N$上で連続である。よって$F$は$X$で連続である。
\end{proof}

\begin{lem}\label{unit}
位相空間$X$とその開被覆$\mathcal{U}$について以下は同値
\begin{align}
&\mbox{$\mathcal{U}$に従属する弱い単位の分割が存在する}\\
&\mbox{$\mathcal{U}$に従属する局所有限な弱い単位の分割が存在する}
\end{align}
\end{lem}
\begin{proof}
$\{f_{\la}\}_{\lil}$
先と同様に
\[
F=\sup_{\lil}f_{\la}
\]
と置く。
そして
\[
g_{\la}=\max\{f_{\la}-\frac{1}{2}F,\ 0\}
\]
と定義する。このとき$\{\wsp(g_{\la})\}_{\lil}$は$X$の被覆となる。これが局所有限であることを示そう。$p\in X$を与えると$F(p)>0$なので
\[
\frac{1}{2}F(p)>\ep
\]
となる$\ep>0$が存在する。そして有限集合$A$で
\[
\sum_{\la \in A}f_{\la}(p)>1-\ep
\]
を充たすものとする。さて
\[
N=\{\sum_{\la \in A}f_{\la}>1-\ep\}\cap \{\frac{1}{2}F>\ep\}
\]
と定義する。すると$N$は開集合で$p\in N$である。ところで
\[
N\cap \wsp(g_{\la})\neq \O
\]
とする。$x\in N\cap \wsp(g_{\la})$を取ると$g_{\la}(x)>0$なので$f_{\la}(x)>\frac{1}{2}F(x)>\ep$となる。よって$f_{\la}(x)>\ep$なので$\la \in A$となるから$N\cap \wsp(g_{\la})\neq \O$となる$\la$は有限個しかない。よって$\{\wsp(g_{\la})\}_{\lil}$は局所有限である。よって
\[
G=\sum_{\la\in \La}g_{\la}
\]
が定義でき、これは連続である。さて
\[
h_{\la}=\frac{g_{\la}}{G}
\]
と置けば
$\wsp(h_{\la})=\wsp(g_{\la})\subset \wsp(f_{\la})$なので確かに$\{h_{\la}\}_{\lil}$は$\mathcal{U}$に従属する局所有限な弱い単位の分割となる。
\end{proof}

\setcounter{equation}{0}
\begin{fact}\label{fact}
位相空間$X$について以下は同値である。
\begin{align}
&\mbox{正規性}\\ 
&\mbox{点有限開被覆に開収縮が存在する}\\ 
&\mbox{局所有限開被覆に開収縮が存在する}\\ 
&\mbox{有限開被覆に開収縮が存在する}\\ 
&\mbox{点有限開被覆に閉収縮が存在する}\\
&\mbox{局所有限開被覆に閉収縮が存在する}\\  
&\mbox{有限開被覆に閉収縮が存在する}
\end{align}
\end{fact}
以下の命題は命題\ref{unit}と事実\ref{fact}から明らかであろう。
\setcounter{equation}{0}
\begin{thm}[局所有限被覆と単位の分割]
位相空間$X$について以下は同値である。
\begin{align}
&\mbox{正規性}\\
&\mbox{任意の局所有限開被覆に対してそれに従属する弱い単位の分割が存在する}\\ 
&\mbox{任意の局所有限開被覆に対してそれに従属する強い単位の分割が存在する}\\ 
&\mbox{任意の局所有限開被覆に対してそれに従属する点有限な弱い単位の分割が存在する}\\ 
&\mbox{任意の局所有限開被覆に対してそれに従属する点有限な強い単位の分割が存在する}\\ 
&\mbox{任意の局所有限開被覆に対してそれに従属する局所有限な弱い単位の分割が存在する}\\ 
&\mbox{任意の局所有限開被覆に対してそれに従属する局所有限な強い単位の分割が存在する}
\end{align}
\end{thm}
\setcounter{equation}{0}
\begin{thm}[有限被覆と単位の分割]
位相空間$X$について以下は同値である。
\begin{align}
&\mbox{正規性}\\
&\mbox{任意の有限開被覆に対してそれに従属する弱い単位の分割が存在する}\\ 
&\mbox{任意の有限開被覆に対してそれに従属する強い単位の分割が存在する}\\ 
&\mbox{任意の有限開被覆に対してそれに従属する点有限な弱い単位の分割が存在する}\\ 
&\mbox{任意の有限開被覆に対してそれに従属する点有限な強い単位の分割が存在する}\\ 
&\mbox{任意の有限開被覆に対してそれに従属する局所有限な弱い単位の分割が存在する}\\ 
&\mbox{任意の有限開被覆に対してそれに従属する局所有限な強い単位の分割が存在する}
\end{align}
\end{thm}






































	
\begin{thebibliography}{9}
\bibitem{1}森田紀一,位相空間論,岩波全書331,岩波書店,1981
\bibitem{2}J.Dydak,\ \textit{Partitions of unity},\ Topology Proceedings 27\ (2003)\ pp125-171
\bibitem{3}
J. Derwent, 
\textit{A note on numerable covers}, 
Proc. Amer. Math. Soc., 
19 (5)(1958), 
1130--1132. 

\end{thebibliography}




			\end{document}



\begin{comment}
[可換図式]
\[\xymatrix{
&   & (2) \ar@{=>}[rd]&  &  &  & (5) \ar@{=>}[rd]&\\ 
& (1)\ar@{=>}[ru] \ar@{=>}[rd]&  &(4)\ar@{=>}[r]&(7)& (1)\ar@{=>}[ru] \ar@{=>}[rd]& & (7)&\\ 
&   & (3) \ar@{=>}[ur]&  &  &  &(6) \ar@{=>}[ur]&
}\]

[\bigin \endを使わない方法]

{\prop[aa] aaa}
で
\begin{prop}[aa]
aaa
\end{prop}
と同じ\df \thm \proof でも同様

[直和]
$\amalg$

[同値関係でよく使う波線]
\sim

[引数付きの新命令]
\newcommand{\prn}[1]{\|#1\|_p}

[番号のリセット]

\setcounter{equation}{0}


[字体の変更]

{\rm hogehoge}と使う。

\rm ローマン
\it イタリック
\sl 斜体

[supp]

\newcommand{\supp}{\text{{\rm supp}}}

[中央揃えの改行]

	\begin{eqnarray*}
		hoge1&=&hogehoge
		hoge2&=&hogehofe

	\end{eqnarray*}

&&の部分で揃えられる。


[\tag]
\begin{equation}
f(x) = ax^2 + bx + c \tag{1.2.3}
\end{equation}



[場合分け]

	\begin{cases}
		hoge1 & \text{状況１}\\
		hoge2 & \text{状況２}\\
		・
		・
		・
	\end{cases}



\end{comment}





